\section{Colombia's payments setting}

\subsection{Migration, documentation and remittance receipt}
Colombia is both a destination for migrants and a country whose households receive income from people living abroad. The central bank's balance-of-payments report puts inward workers' remittances at US\$13.1 billion in calendar year 2025, equivalent to 2.9\% of national GDP \citep{banrepbop2026}. Those are national flows into Colombia. They do not describe the Colombian--Venezuelan corridor in both directions, remittances sent by women living in Colombia, or the frequency and amount received by our survey respondents. Recipient households, individual users, and aggregate cross-border flows are different units.

The documentation environment changed before fieldwork. The PEP evaluated in earlier regularization research and the \emph{Permiso por Protección Temporal} (PPT) are distinct. Decree 216 of 2021 established a temporary protection statute for eligible Venezuelan migrants and created the PPT and the associated register \citep{decree216}. \href{https://www.cancilleria.gov.co/normograma/compilacion/docs/pdf/resolucion_uaemc_0971_2021.pdf}{Article 14 of Migration Resolution 971 of 2021} recognizes the PPT as identification for products and services at supervised financial institutions, subject to other applicable requirements. The PPT is not a citizenship credential. The survey does not verify legal status, product eligibility or whether a provider accepted a document. Using someone else's account can solve an immediate problem without giving the user continuing access or control.

\subsection{From separate services to interoperable payments}
Colombia's payment market predates Bre-B. Law 1735 of 2014 created the framework for specialized electronic deposits and payments companies, extending the institutional options for offering payment and deposit services \citep{law1735}. Wallets and private instant-transfer initiatives developed before national interoperation. A payment possible within one service, or through a bilateral arrangement, did not necessarily reach customers of other providers. Article 104 of Law 2294, External Resolution 6 of 2023 and Circular DSP-465 establish the regulatory basis for interoperability among low-value immediate payment systems \citep{banrepreg,banreptechnical2026}. This framework set rules for interoperation before providers began full operation.

Alias registration began in July 2025. Controlled operation ran from 23 September through 5 October; full operation began on 6 October 2025 \citep{banreppilot,banreptechnical2026,banrep2026}. Bre-B is not a standalone app. Participating providers offer the service through existing channels; aliases in the central directory and the central settlement arrangement serve different functions. The user who sends a transfer need not register an alias solely to send; receiving through an alias depends on the recipient's linked account or deposit. For the survey period, the documented Bre-B scheme enabled domestic payments and transfers \citep{banrepfaq,banreptechnical2026}. That scope does not establish the origin, conversion or settlement of an incoming international remittance. A wallet can receive funds through a separate arrangement and then use a domestic rail for a subsequent payment.

Registering for the service and making a payment are different actions. The technical report distinguishes registered keys, customers, participating entities and operations settled on the new mechanism \citep{banreptechnical2026}. A customer can register more than one key, and registration does not establish use. Aggregate payment volume can rise without indicating which population gained access, benefited from lower prices, or had a complaint resolved. The survey does not identify the processing rail or the date of first use. The recorded survey interviews, from 9 December 2025 through 17 February 2026, came after full operation; some item reference periods extend farther back than others. Neither survey dates nor recall periods create an exposure measure or a before-and-after design.

\subsection{The conditions around the rail}
Connectivity provides one reason to distinguish technical availability from access. DANE's 2024 Quality of Life Survey reports internet connectivity for 65.6\% of households nationally, 72.5\% in municipal seats and 41.9\% in populated centers and dispersed rural areas \citep{daneentic2025}. These are household estimates for a pre-launch reference year, not the proportion of migrant women who have their own connected device. Global Findex 2025, by contrast, uses interviews with adults conducted in 2024 \citep{findex2025}. Household connectivity, adults' reported financial behavior and provider-reported registrations should not be joined into one adoption rate. For an individual customer, changing a phone number, lacking sufficient data, or needing help at a verification step can still interrupt use in a connected area.

Rules for sharing financial data are also distinct from rules for sending funds. The 2022 open-finance framework preceded the fieldwork; Decree 0368, issued in April 2026, set a later mandatory framework \citep{decree1297,decree0368}. It is not part of the policy environment during the recorded survey interviews, and promulgation is not evidence that every data-sharing service was operating. In domestic Bre-B transactions, the central bank directs customers to the financial entities involved to pursue disputes \citep{banrepfaq}. This documents a route, not its accessibility in practice, a guaranteed refund, or responsibility for a foreign settlement or an intermediary's actions. Public oversight of a network and a user's ability to recover funds are separate questions.

Identity rules and payment-system design therefore answer different questions. One concerns who can open and keep a usable account; the other concerns which domestic accounts can exchange payments. Neither records what an international recipient receives or what she does with the money.
